\subsection{Separated affinoids for the strong maximum principle}\label{subsec:fc-line-separated-neighbourhoods}

\begin{definition}[Separated presentation and strict core]\label{def:fc-separated-affinoid-core}
\lean{EnvelopeProperty.Berkovich.RelativeSpectrum.FCStandardAffinoid.IsSeparated, EnvelopeProperty.Berkovich.RelativeSpectrum.FCStandardAffinoid.strictCore}
\leanok
\uses{def:fc-standard-affinoid}
A presentation of an actual standard affinoid is separated if each hole
centre is strictly inside the outer radius and, for distinct holes,
each hole radius is strictly smaller than their centre distance.
Its strict core is the set of actual points satisfying all its defining
inequalities strictly.
\end{definition}

\begin{theorem}[Separated actual neighbourhoods]\label{thm:fc-separated-affinoid-basis}
\lean{EnvelopeProperty.Berkovich.RelativeSpectrum.fc_exists_separated_standardAffinoid_nhds}
\leanok
\uses{def:fc-separated-affinoid-core, thm:fc-standard-affinoid-basis,
  lem:fc-linear-topology}
Every neighbourhood $N$ of every actual line point $x$ contains the closed
carrier of an actual separated standard affinoid whose strict core contains
$x$. This is a construction, not an additional hypothesis on the source
definition of a standard affinoid.
\end{theorem}

\begin{proof}
Choose a finite box in original linear coordinates inside $N$. Between the
outer coordinate value at $x$ and the smallest permitted upper bound choose
an outer radius avoiding the finitely many centre distances. For each
positive lower bound choose a slightly larger radius still strictly below
that coordinate value at $x$, avoiding all the finitely many distances
between box centres. Keep zero lower bounds zero. Select maximal members
of the resulting finite laminar open-disc family. They are disjoint and
cover all the required excluded discs. The chosen avoidance inequalities
turn every weak centre-distance bound into a strict one. The original
point satisfies every resulting defining inequality strictly, and the
closed carrier remains in the original coordinate box.
\end{proof}

\begin{lemma}[The connected core approaches every actual boundary]\label{lem:fc-separated-core-connected-boundary}
\lean{EnvelopeProperty.Berkovich.RelativeSpectrum.FCStandardAffinoid.isPathConnected_strictCore, EnvelopeProperty.Berkovich.RelativeSpectrum.FCStandardAffinoid.boundary_subset_closure_strictCore}
\leanok
\uses{def:fc-separated-affinoid-core, lem:fc-inflation-linear}
The nonempty strict core of an actual separated standard affinoid is
path connected. Every actual Gauss boundary point lies in its closure.
\end{lemma}

\begin{proof}
Two strict-core points inflate to the same outer-centred Gauss point at a
radius smaller than the outer radius and larger than both outer coordinate
values and all hole radii and centre distances. Inflation increases all
linear coordinates, so both paths remain in the strict core.
An inner boundary is approached by Gauss points of the same centre at
radii just above its hole radius and below all other centre distances and
the outer radius. The outer boundary is approached from radii just below
the outer radius and above all the finitely many hole radii and centre
distances. Coincident boundaries and zero-length segments need no
assumptions here: this separated presentation is constructed solely for
this argument and is not imposed on arbitrary standard affinoids.
\end{proof}
